\documentclass[../../main.tex]{subfiles}
\begin{document}
\section{子空间、积、不交并与商}\label{section:appA-subspaces-products-quotients}

\subsection*{子空间}

从旧空间获得新拓扑空间最简单的办法，或许是取其他空间的子集。如果 $X$ 是拓扑空间而 $S\subset X$ 是任意子集，我们在 $S$ 上定义\emph{子空间拓扑}（有时也称相对拓扑）：宣告子集 $U\subset S$ 在 $S$ 中开，当且仅当存在开子集 $V\subset X$ 使 $U=V\cap S$。在子空间拓扑中开或闭的 $S$ 的子集，有时称为 $S$ 中的\emph{相对开}或\emph{相对闭}子集，以表明我们不是指作为 $X$ 子集的开或闭。赋予子空间拓扑的 $X$ 的任意子集称为 $X$ 的一个\emph{子空间}。每当我们把拓扑空间的一个子集本身当作空间处理时，除非另行说明，总假定它带子空间拓扑。

如果 $X$ 与 $Y$ 是拓扑空间，连续单射 $F:X\to Y$ 如果是到其像 $F(X)\subset Y$（带子空间拓扑）的同胚，就称为\emph{拓扑嵌入}。

子空间拓扑最重要的性质概括在下一个命题中。

\begin{proposition}[子空间拓扑的性质]\label{proposition:appA-subspace-properties}
设 $X$ 是拓扑空间，$S$ 是 $X$ 的子空间。
\begin{enumerate}[(a)]
\item \textbf{特征性质：}如果 $Y$ 是拓扑空间，映射 $F:Y\to S$ 连续，当且仅当复合 $\iota_S\circ F:Y\to X$ 连续，其中 $\iota_S:S\hookrightarrow X$ 是包含映射（$X$ 的恒等映射到 $S$ 的限制）。
\item 子空间拓扑是 $S$ 上使特征性质成立的唯一拓扑。
\item 子集 $K\subset S$ 在 $S$ 中闭，当且仅当存在闭子集 $L\subset X$ 使 $K=L\cap S$。
\item 包含映射 $\iota_S:S\hookrightarrow X$ 是拓扑嵌入。
\item 如果 $Y$ 是拓扑空间而 $F:X\to Y$ 连续，则 $F|_S:S\to Y$（$F$ 到 $S$ 的限制）连续。
\item 如果 $\mathcal B$ 是 $X$ 拓扑的一个基，则 $\mathcal B_S=\{B\cap S:B\in\mathcal B\}$ 是 $S$ 上子空间拓扑的一个基。
\item 如果 $X$ 是 Hausdorff 的，则 $S$ 也是。
\item 如果 $X$ 是第一可数的，则 $S$ 也是。
\item 如果 $X$ 是第二可数的，则 $S$ 也是。
\end{enumerate}
\end{proposition}

\begin{exercise}\label{exercise:appA-A-18}
证明前一个命题。
\end{exercise}

如果 $X$ 与 $Y$ 是拓扑空间而 $F:X\to Y$ 连续，前一命题的 (e) 保证 $F$ 到 $X$ 的每个子空间的限制（在子空间拓扑下）都连续。我们也可以问逆问题：如果已知 $F$ 到 $X$ 的某些子空间的限制连续，$F$ 本身是否连续？下列两个命题给出对这个问题两个略有不同的回答。

\begin{lemma}[连续性是局部性质]\label{lemma:appA-continuity-local}
连续性是局部性质，其意义如下：如果 $F:X\to Y$ 是拓扑空间之间的映射，且每个点 $p\in X$ 都有邻域 $U$ 使限制 $F|_U$ 连续，则 $F$ 连续。
\end{lemma}

\begin{lemma}[连续映射的粘合引理]\label{lemma:appA-gluing-lemma}
设 $X$ 与 $Y$ 是拓扑空间，并假设下列条件之一成立：
\begin{enumerate}[(a)]
\item $B_1,\dots,B_n$ 是 $X$ 的有限个闭子集，其并为 $X$。
\item $\{B_i\}_{i\in A}$ 是 $X$ 的一族开子集，其并为 $X$。
\end{enumerate}
假设对所有 $i$，我们给定连续映射 $F_i:B_i\to Y$，且它们在重叠处一致：$F_i|_{B_i\cap B_j}=F_j|_{B_i\cap B_j}$。则存在唯一连续映射 $F:X\to Y$，其到每个 $B_i$ 的限制等于 $F_i$。
\end{lemma}

\begin{exercise}\label{exercise:appA-A-21}
证明前两个引理。
\end{exercise}

\begin{exercise}\label{exercise:appA-A-22}
设 $X$ 是拓扑空间，并假设 $X$ 有可数开覆盖 $\{U_i\}$，使每个集合 $U_i$ 在子空间拓扑下都是第二可数的。证明 $X$ 是第二可数的。
\end{exercise}

\subsection*{积空间}

接下来我们考虑拓扑空间的有限积。如果 $X_1,\dots,X_k$ 是（有限个）集合，它们的\emph{Cartesian 积}是集合 $X_1\times\cdots\times X_k$，由所有形如 $(x_1,\dots,x_k)$（其中对每个 $i$ 都有 $x_i\in X_i$）的有序 $k$ 元组组成。第 $i$ 个\emph{投影映射}是映射 $\pi_i:X_1\times\cdots\times X_k\to X_i$，定义为 $\pi_i(x_1,\dots,x_k)=x_i$。

假设 $X_1,\dots,X_k$ 是拓扑空间。形如 $U_1\times\cdots\times U_k$（其中每个 $U_i$ 在 $X_i$ 中开）的 $X_1\times\cdots\times X_k$ 的所有子集所成族，构成 $X_1\times\cdots\times X_k$ 上一个拓扑的基，称为\emph{积拓扑}。赋予此拓扑，有限个拓扑空间的积称为一个\emph{积空间}。形如 $U_1\times\cdots\times U_k\subset X_1\times\cdots\times X_k$（其中每个 $U_i$ 在 $X_i$ 中开）的任意开子集称为\emph{积开子集}。（无穷多个空间的积需要略为不同的定义，但我们只需要有限情形。关于无穷积空间见 [LeeTM]。）

\begin{proposition}[积拓扑的性质]\label{proposition:appA-product-properties}
假设 $X_1,\dots,X_k$ 是拓扑空间，$X_1\times\cdots\times X_k$ 是它们的积空间。
\begin{enumerate}[(a)]
\item \textbf{特征性质：}如果 $B$ 是拓扑空间，映射 $F:B\to X_1\times\cdots\times X_k$ 连续，当且仅当其每个分量函数 $F_i=\pi_i\circ F:B\to X_i$ 都连续。
\item 积拓扑是 $X_1\times\cdots\times X_k$ 上使特征性质成立的唯一拓扑。
\item 每个投影映射 $\pi_i:X_1\times\cdots\times X_k\to X_i$ 连续。
\item 给定任意连续映射 $F_i:X_i\to Y_i$（$i=1,\dots,k$），积映射 $F_1\times\cdots\times F_k:X_1\times\cdots\times X_k\to Y_1\times\cdots\times Y_k$ 连续，其中
\begin{align*}
F_1\times\cdots\times F_k(x_1,\dots,x_k)=(F_1(x_1),\dots,F_k(x_k)).
\end{align*}
\item 如果 $S_i$ 是 $X_i$（$i=1,\dots,n$）的子空间，则 $S_1\times\cdots\times S_n\subset X_1\times\cdots\times X_n$ 上的积拓扑与子空间拓扑一致。
\item 对任意 $i\in\{1,\dots,k\}$ 和任意 $j\neq i$ 处点 $a_j\in X_j$ 的选取，映射 $x\mapsto(a_1,\dots,a_{i-1},x,a_{i+1},\dots,a_k)$ 是 $X_i$ 到积空间 $X_1\times\cdots\times X_k$ 的拓扑嵌入。
\item 如果 $\mathcal B_i$ 是 $X_i$（$i=1,\dots,k$）拓扑的基，则族 $\mathcal B=\{B_1\times\cdots\times B_k:B_i\in\mathcal B_i\}$ 是 $X_1\times\cdots\times X_k$ 拓扑的基。
\item Hausdorff 空间的任意有限积是 Hausdorff 的。
\item 第一可数空间的任意有限积是第一可数的。
\item 第二可数空间的任意有限积是第二可数的。
\end{enumerate}
\end{proposition}

\begin{exercise}\label{exercise:appA-A-24}
证明前一个命题。
\end{exercise}

\subsection*{不交并空间}

构造新拓扑空间的另一个简单办法是取其他空间的不交并。从集合论观点看，不交并定义如下。如果 $(X_\alpha)_{\alpha\in A}$ 是加标集合族，它们的\emph{不交并}是集合
\begin{align*}
\coprod_{\alpha\in A}X_\alpha=\{(x,\alpha):\alpha\in A,\ x\in X_\alpha\}.
\end{align*}
对每个 $\alpha$，有典范单射 $\iota_\alpha:X_\alpha\to\coprod_{\alpha\in A}X_\alpha$，由 $\iota_\alpha(x)=(x,\alpha)$ 给出，而不同 $\alpha$ 处这些映射的像不相交。通常我们默认把 $X_\alpha$ 与其在不交并中的像等同起来，从而把 $X_\alpha$ 看作 $\coprod_{\alpha\in A}X_\alpha$ 的子集。记号 $(x,\alpha)$ 中的 $\alpha$ 应看作一个“标签”，用以表明 $x$ 来自哪个集合，使不同 $\alpha$ 对应的子集不相交，即使某些或全部原集合 $X_\alpha$ 相同。

给定加标拓扑空间族 $(X_\alpha)_{\alpha\in A}$，我们在 $\coprod_{\alpha\in A}X_\alpha$ 上定义\emph{不交并拓扑}：宣告 $\coprod_{\alpha\in A}X_\alpha$ 的一个子集是开的，当且仅当它与每个 $X_\alpha$ 的交在 $X_\alpha$ 中开。

\begin{proposition}[不交并拓扑的性质]\label{proposition:appA-disjoint-union-properties}
假设 $(X_\alpha)_{\alpha\in A}$ 是加标拓扑空间族，$\coprod_{\alpha\in A}X_\alpha$ 赋予不交并拓扑。
\begin{enumerate}[(a)]
\item \textbf{特征性质：}如果 $Y$ 是拓扑空间，映射 $F:\coprod_{\alpha\in A}X_\alpha\to Y$ 连续，当且仅当对每个 $\alpha\in A$，$F\circ\iota_\alpha:X_\alpha\to Y$ 连续。
\item 不交并拓扑是 $\coprod_{\alpha\in A}X_\alpha$ 上使特征性质成立的唯一拓扑。
\item $\coprod_{\alpha\in A}X_\alpha$ 的一个子集闭，当且仅当它与每个 $X_\alpha$ 的交都闭。
\item 每个单射 $\iota_\alpha:X_\alpha\to\coprod_{\alpha\in A}X_\alpha$ 是拓扑嵌入。
\item Hausdorff 空间的任意不交并是 Hausdorff 的。
\item 第一可数空间的任意不交并是第一可数的。
\item 可数个第二可数空间的不交并是第二可数的。
\end{enumerate}
\end{proposition}

\begin{exercise}\label{exercise:appA-A-26}
证明前一个命题。
\end{exercise}

\subsection*{商空间与商映射}

如果 $X$ 是拓扑空间，$Y$ 是集合，$\pi:X\to Y$ 是满射，则由 $\pi$ 确定的 $Y$ 上\emph{商拓扑}定义为：宣告子集 $U\subset Y$ 是开的，当且仅当 $\pi^{-1}(U)$ 在 $X$ 中开。如果 $X$ 与 $Y$ 是拓扑空间，映射 $\pi:X\to Y$ 如果是满射、连续且 $Y$ 带有由 $\pi$ 确定的商拓扑，就称为\emph{商映射}。

下列构造是产生商映射最常见的方式。集合 $X$ 上的一个关系 $\sim$ 如果是自反的（对所有 $x\in X$ 有 $x\sim x$）、对称的（$x\sim y$ 蕴含 $y\sim x$）、传递的（$x\sim y$ 且 $y\sim z$ 蕴含 $x\sim z$），就称为一个\emph{等价关系}。如果 $R\subset X\times X$ 是 $X$ 上任意关系，则 $X$ 上所有含 $R$ 的等价关系之交是一个等价关系，称为由 $R$ \emph{生成的等价关系}。如果 $\sim$ 是 $X$ 上的等价关系，则对每个 $x\in X$，$x$ 的\emph{等价类}记作 $[x]$，是所有满足 $y\sim x$ 的 $y\in X$ 的集合。所有等价类所成集合是 $X$ 的一个\emph{划分}：一族不相交的非空子集，其并为 $X$。

假设 $X$ 是拓扑空间而 $\sim$ 是 $X$ 上的等价关系。令 $X/{\sim}$ 表示 $X$ 中等价类所成集合，$\pi:X\to X/{\sim}$ 是把每点映到其等价类的自然投影。赋予由 $\pi$ 确定的商拓扑，空间 $X/{\sim}$ 称为由 $\pi$ 确定的 $X$ 的\emph{商空间}（或\emph{粘合空间}）。例如，假设 $X$ 与 $Y$ 是拓扑空间，$A\subset Y$ 是闭子集，$f:A\to X$ 是连续映射。关系“对所有 $a\in A$ 有 $a\sim f(a)$”在 $X\sqcup Y$ 上生成一个等价关系，其商空间记作 $X\cup_f Y$，称为一个\emph{贴附空间}。它被说成是“沿 $f$ 把 $Y$ 贴附到 $X$”而形成的。

如果 $\pi:X\to Y$ 是映射，子集 $U\subset X$ 如果是其像的整个原像：$U=\pi^{-1}(\pi(U))$，就称为关于 $\pi$ \emph{饱和的}。给定 $y\in Y$，$\pi$ 在 $y$ 上的\emph{纤维}是集合 $\pi^{-1}(y)$。因此，$X$ 的子集饱和，当且仅当它是纤维之并。

\begin{theorem}[商映射的性质]\label{theorem:appA-quotient-properties}
设 $\pi:X\to Y$ 是商映射。
\begin{enumerate}[(a)]
\item \textbf{特征性质：}如果 $B$ 是拓扑空间，映射 $F:Y\to B$ 连续，当且仅当 $F\circ\pi:X\to B$ 连续。
\item 商拓扑是 $Y$ 上使特征性质成立的唯一拓扑。
\item 子集 $K\subset Y$ 闭，当且仅当 $\pi^{-1}(K)$ 在 $X$ 中闭。
\item 如果 $\pi$ 是单射，则它是同胚。
\item 如果 $U\subset X$ 是饱和的开或闭子集，则限制 $\pi|_U:U\to\pi(U)$ 是商映射。
\item $\pi$ 与另一个商映射的任意复合仍是商映射。
\end{enumerate}
\end{theorem}

\begin{exercise}\label{exercise:appA-A-28}
证明前一个定理。
\end{exercise}

\begin{exercise}\label{exercise:appA-A-29}
设 $X$ 与 $Y$ 是拓扑空间，假设 $F:X\to Y$ 是满连续映射。证明下列各条等价：
\begin{enumerate}[(a)]
\item $F$ 是商映射。
\item $F$ 把饱和开子集映为开子集。
\item $F$ 把饱和闭子集映为闭子集。
\end{enumerate}
\end{exercise}

商映射的下列两个性质在拓扑学中起重要作用，并在光滑流形理论中有同等重要的推广（见第 4 章）。

\begin{theorem}[过渡到商]\label{theorem:appA-passing-to-quotient}
假设 $\pi:X\to Y$ 是商映射，$B$ 是拓扑空间，$F:X\to B$ 是连续映射且在 $\pi$ 的纤维上为常数（即 $\pi(p)=\pi(q)$ 蕴含 $F(p)=F(q)$）。则存在唯一连续映射 $\widetilde F:Y\to B$ 使 $F=\widetilde F\circ\pi$。
\end{theorem}

\begin{proof}
$\widetilde F$ 的存在与唯一性由集合论考虑得到，而其连续性是商拓扑特征性质的直接推论。
\end{proof}

\begin{theorem}[商空间的唯一性]\label{theorem:appA-uniqueness-quotient}
如果 $\pi_1:X\to Y_1$ 与 $\pi_2:X\to Y_2$ 是在彼此纤维上为常数的商映射（即 $\pi_1(p)=\pi_1(q)$ 当且仅当 $\pi_2(p)=\pi_2(q)$），则存在唯一同胚 $\varphi:Y_1\to Y_2$ 使 $\varphi\circ\pi_1=\pi_2$。
\end{theorem}

\begin{proof}
把前一定理应用于商映射 $\pi_1:X\to Y_1$，我们看到 $\pi_2$ 过渡到商，得到连续映射 $\widetilde\pi_2:Y_1\to Y_2$ 满足 $\widetilde\pi_2\circ\pi_1=\pi_2$。把同样论证中 $\pi_1$ 与 $\pi_2$ 的角色互换，有连续映射 $\widetilde\pi_1:Y_2\to Y_1$ 满足 $\widetilde\pi_1\circ\pi_2=\pi_1$。合起来，这些恒等式蕴含 $\widetilde\pi_2\circ\widetilde\pi_1\circ\pi_2=\pi_2$。再把定理 A.30 应用于让 $\pi_2$ 同时扮演 $\pi$ 与 $F$ 的角色，我们看到 $\widetilde\pi_2\circ\widetilde\pi_1$ 与 $\operatorname{Id}_{Y_2}$ 都由 $\pi_2$ 通过过渡到商得到，故定理 A.30 的唯一性断言蕴含 $\widetilde\pi_2\circ\widetilde\pi_1=\operatorname{Id}_{Y_2}$。类似论证表明 $\widetilde\pi_1\circ\widetilde\pi_2=\operatorname{Id}_{Y_1}$，故 $\widetilde\pi_2$ 就是所求同胚。
\end{proof}

\subsection*{开映射与闭映射}

映射 $F:X\to Y$（无论是否连续）如果对每个开子集 $U\subset X$，像集 $F(U)$ 在 $Y$ 中开，就称为\emph{开映射}；如果对每个闭子集 $K\subset X$，像 $F(K)$ 在 $Y$ 中闭，就称为\emph{闭映射}。连续映射可以是开的、闭的、既开又闭的，或既不开也不闭，这可由考察涉及平面子集的简单例子看出。

\begin{exercise}\label{exercise:appA-A-32}
假设 $X_1,\dots,X_k$ 是拓扑空间。证明每个投影 $\pi_i:X_1\times\cdots\times X_k\to X_i$ 是开映射。
\end{exercise}

\begin{exercise}\label{exercise:appA-A-33}
设 $(X_\alpha)_{\alpha\in A}$ 是加标拓扑空间族。证明每个单射 $\iota_\alpha:X_\alpha\to\coprod_{\alpha\in A}X_\alpha$ 既开又闭。
\end{exercise}

\begin{exercise}\label{exercise:appA-A-34}
证明每个局部同胚是开映射。
\end{exercise}

\begin{exercise}\label{exercise:appA-A-35}
证明每个双射局部同胚是同胚。
\end{exercise}

\begin{exercise}\label{exercise:appA-A-36}
假设 $q:X\to Y$ 是开商映射。证明 $Y$ 是 Hausdorff 的，当且仅当集合 $R=\{(x_1,x_2):q(x_1)=q(x_2)\}$ 在 $X\times X$ 中闭。
\end{exercise}

\begin{exercise}\label{exercise:appA-A-37}
设 $X$ 与 $Y$ 是拓扑空间，$F:X\to Y$ 是映射。证明下列各条：
\begin{enumerate}[(a)]
\item $F$ 闭，当且仅当对每个 $A\subset X$，$F(\bar A)\supset\overline{F(A)}$。
\item $F$ 开，当且仅当对每个 $B\subset Y$，$F^{-1}(\operatorname{Int}B)\supset\operatorname{Int}F^{-1}(B)$。
\end{enumerate}
\end{exercise}

拓扑学中最重要的连续映射类是同胚、商映射和拓扑嵌入。显然，一个映射要成为同胚必须是双射，成为商映射必须是满射，成为拓扑嵌入必须是单射。然而，即使已知连续映射满足这些必要的集合论条件之一，也并不总能容易地判断它是否有所需的拓扑性质。如下一个定理所示，一个简单的充分条件是它为开映射或闭映射。

\begin{theorem}\label{theorem:appA-open-closed-maps}
假设 $X$ 与 $Y$ 是拓扑空间，$F:X\to Y$ 是连续映射且为开或闭。
\begin{enumerate}[(a)]
\item 如果 $F$ 是满射，则它是商映射。
\item 如果 $F$ 是单射，则它是拓扑嵌入。
\item 如果 $F$ 是双射，则它是同胚。
\end{enumerate}
\end{theorem}

\begin{proof}
先设 $F$ 是满射。如果它是开的，它当然把饱和开子集映为开子集。类似地，如果它是闭的，它把饱和闭子集映为闭子集。故由习题 A.29，它是商映射。

现设 $F$ 是开且单射。则 $F:X\to F(X)$ 是双射，故由初等集合论考虑 $F^{-1}:F(X)\to X$ 存在。如果 $U\subset X$ 开，则 $(F^{-1})^{-1}(U)=F(U)$ 由假设在 $Y$ 中开，因而由 $F(X)$ 上子空间拓扑的定义在 $F(X)$ 中也开。这证明 $F^{-1}$ 连续，故 $F$ 是到其像的同胚。如果 $F$ 是闭的，同样的论证把“开”换成“闭”即可（使用命题 A.17(c) 中给出的 $F(X)$ 闭子集的刻画）。这就证明了 (b)，而 (c) 只是 (b) 中 $F(X)=Y$ 的特殊情形。
\end{proof}
\end{document}
